% =====================================================================
% ejercicio05.tex
% Responsabilidad única: contrastar dos matrices de confusión bajo las
% cuatro métricas estándar y bajo un costo asimétrico L=FP+5FN,
% evidenciando que "mejor Accuracy" no implica "menor costo".
% =====================================================================
\subsection{Ejercicio 5 --- Comparación de dos modelos y costo de clasificación}

\begin{table}[H]
\centering
\begin{minipage}{0.45\textwidth}
\centering
\begin{tabular}{cc|cc}
    & & \multicolumn{2}{c}{\textbf{Predicho}} \\
    & & $\hat Y=0$ & $\hat Y=1$ \\
    \hline
    \multirow{2}{*}{\textbf{Real}} & $Y=0$ & 900 & 50 \\
    & $Y=1$ & 30 & 20 \\
\end{tabular}
\caption*{Modelo A}
\end{minipage}
\hfill
\begin{minipage}{0.45\textwidth}
\centering
\begin{tabular}{cc|cc}
    & & \multicolumn{2}{c}{\textbf{Predicho}} \\
    & & $\hat Y=0$ & $\hat Y=1$ \\
    \hline
    \multirow{2}{*}{\textbf{Real}} & $Y=0$ & 850 & 100 \\
    & $Y=1$ & 10 & 40 \\
\end{tabular}
\caption*{Modelo B}
\end{minipage}
\caption{Matrices de confusión de los dos modelos comparados.}
\end{table}

\begin{enumerate}

\item \textbf{Métricas de cada modelo.}
Modelo A ($TN=900,FP=50,FN=30,TP=20$):
\[
    \text{Acc}_A = 0.92, \quad
    \text{Prec}_A = \tfrac{20}{70}\approx 0.2857, \quad
    \text{Rec}_A = 0.4, \quad
    F1_A = \tfrac13 \approx 0.3333.
\]
Modelo B ($TN=850,FP=100,FN=10,TP=40$):
\[
    \text{Acc}_B = 0.89, \quad
    \text{Prec}_B = \tfrac{40}{140}\approx 0.2857, \quad
    \text{Rec}_B = 0.8, \quad
    F1_B = \tfrac{8}{19} \approx 0.4211.
\]

\item \textbf{Costo $L=FP+5FN$.}
\[
    L_A = 50+5(30) = 200, \qquad L_B = 100+5(10) = 150.
\]

\item \textbf{Mayor Accuracy:} el Modelo A ($0.92 > 0.89$).

\item \textbf{Mayor Recall:} el Modelo B ($0.8 > 0.4$), por amplio margen.

\item \textbf{Menor costo:} el Modelo B ($150 < 200$).

\item \textbf{El modelo de mayor Accuracy no minimiza el costo.}
El Modelo A tiene mayor Accuracy pero un mayor coste asociado que el
Modelo B, esto sucede porque lo que hace el Accuracy es ponderar todos
los errores por igual, mientras que $L$ pondera cada falso negativo
cinco veces más que cada falso positivo: el Modelo A minimiza el
número total de errores ($80$ en vez de $110$), pero comete más falsos
negativos ($30$ y no $10$), que son los más costosos, mientras que el
Modelo B sí compensa con más falsos positivos ($100$ en vez de $50$),
los que son baratos. Cuando los dos tipos de error tienen consecuencias
distintas, elegir un modelo únicamente por el Accuracy puede ser una
decisión subóptima enmarcada en el costo real del negocio.

\end{enumerate}
